\documentclass[aspectratio=169]{beamer}
\usetheme{Madrid}
\usecolortheme{default}
\setbeamertemplate{navigation symbols}{}

\usepackage[T1]{fontenc}
\usepackage[utf8]{inputenc}
\usepackage{lmodern}

\title{ERC-3 EarthRover Project Update}
\subtitle{Indoor, Outdoor, and Marathon Summary}
\author{Lunar}
\date{April 1, 2026}

\begin{document}

\begin{frame}
  \titlepage
\end{frame}

\begin{frame}{Current Status Update: Later No-GPS Branch}
\textbf{This deck is an April 2026 snapshot of indoor, GPS-outdoor, and marathon work.}

\vspace{0.5em}
\begin{itemize}
  \item A later branch uses front-camera-only visual teach-and-repeat because GPS is unavailable.
  \item Recording, route preparation, localization, and live command sending work.
  \item A reliable autonomous physical route repeat is \textbf{not yet proven}.
  \item The correct current claim is: experimental prototype, not competition-ready autonomy.
\end{itemize}

\vspace{0.5em}
\small Current evidence: \texttt{obsidian\_vault/01 Source of Truth/No-GPS Field Trial - Findings.md}
\end{frame}

\begin{frame}{Indoor}
\textbf{What we chose}
\begin{itemize}
  \item MBRA as the main indoor controller
  \item exact checkpoint-step navigation
  \item a simpler runtime around it
\end{itemize}

\vspace{0.5em}
\textbf{Why we chose it}
\begin{itemize}
  \item indoor is a known corridor problem
  \item MBRA worked better with simpler recovery
  \item exact steps were cleaner than vague targets
\end{itemize}
\end{frame}

\begin{frame}{Indoor Results}
\textbf{What worked}
\begin{itemize}
  \item \textbf{8 / 11 checkpoints reached}
  \item the system became more stable
  \item the MBRA-first direction became clearer
\end{itemize}

\vspace{0.5em}
\textbf{What failed / what was difficult}
\begin{itemize}
  \item forward-only graph caused no-path issues
  \item stale context could trap the robot
  \item recovery still needed care
\end{itemize}
\end{frame}

\begin{frame}{Outdoor}
\textbf{What we chose}
\begin{itemize}
  \item LogoNav as the main outdoor controller
  \item OSM route expansion
  \item safety layers around the runtime
\end{itemize}

\vspace{0.5em}
\textbf{Why we chose it}
\begin{itemize}
  \item outdoor was not a from-scratch controller problem
  \item we already had a working base runtime
  \item the bigger need was safety and stability
\end{itemize}
\end{frame}

\begin{frame}{Outdoor Results}
\textbf{What worked}
\begin{itemize}
  \item one run reached \textbf{all checkpoints}
  \item the others were around \textbf{50\% success}
  \item rerouting and waypoint fixes helped a lot
\end{itemize}

\vspace{0.5em}
\textbf{What failed / what was difficult}
\begin{itemize}
  \item \textbf{3 rounds failed}
  \item spinning and bad waypoint transitions hurt us
  \item curbs and steps are still a weak point
\end{itemize}
\end{frame}

\begin{frame}{Marathon}
\textbf{What we chose}
\begin{itemize}
  \item keep the same outdoor base runtime
  \item add stricter safety layers
  \item focus on stability, not full autonomy
\end{itemize}

\vspace{0.5em}
\textbf{What we added}
\begin{itemize}
  \item IMU safety
  \item route corridor guard
  \item rerouting from live GPS
  \item better waypoint handling
  \item clearer logs
\end{itemize}
\end{frame}

\begin{frame}{Marathon Result And Failure Reason}
\Large
\centering
\textbf{Reached 1 checkpoint}

\vspace{0.8em}
\textbf{Then the robot spun and toppled}

\normalsize
\vspace{1.0em}
\textbf{What happened}
\begin{itemize}
  \item after checkpoint progress, the next target handling became unstable
  \item that led to repeated turning in place
\end{itemize}

\vspace{0.5em}
\textbf{Why this likely happened}
\begin{itemize}
  \item the runtime likely got confused after checkpoint transition
  \item that created bad alignment behavior
  \item repeated aggressive turning made the robot unstable
\end{itemize}

\vspace{0.5em}
\textbf{Main lesson}

The marathon problem was not just navigation. It was stability after checkpoint and waypoint transitions.
\end{frame}

\begin{frame}{Final Summary}
\textbf{Indoor}
\begin{itemize}
  \item best current direction: MBRA-first
  \item result: \textbf{8 / 11}
\end{itemize}

\textbf{Outdoor}
\begin{itemize}
  \item best current direction: LogoNav + route + safety layers
  \item result: \textbf{1 full success}, others around \textbf{50\%}
\end{itemize}

\textbf{Marathon}
\begin{itemize}
  \item result: \textbf{1 checkpoint, then toppled}
  \item next focus: stability after transitions
\end{itemize}
\end{frame}

\end{document}
